\section*{Original degree-bounded evasion convention (journal p.6)}
\begin{result}{Definition 1.5}
We say a projective (resp. affine) $k$-subspace family $\mathcal H$ on $\PP^n$ (resp. $\Aspace^n$) is $(n,d)$-evasive if it is $\mathcal F$-evasive, where $\mathcal F$ is chosen to be the family of all subvarieties of $\PP^n$ (resp. $\Aspace^n$) of degree at most $d$. Similarly, we say $\mathcal H$ is $(n,d,\epsilon)$-evasive if it is $(\mathcal F,\epsilon)$-evasive.
\end{result}
\begin{result}{Remark 1.6}
In Definition 1.5, we do not make any assumption about the dimension of the varieties in $\mathcal F$ or their irreducible components. We will see in Section 3.1 that in fact, it suffices to consider the subfamily of equidimensional varieties or even irreducible varieties of dimension $n-k-1$ when constructing variety evasive $k$-subspace families.
\end{result}
\section*{Original Section 2.1: Black-Box PIT for Low Degree Polynomials (journal pp.18--20)}
\begin{result}{Definition 2.1 ($\epsilon$-hitting set)}
Let $\mathcal F$ be a family of polynomials in $\FF[X_1,\ldots,X_n]$ and $\epsilon\in(0,1)$. We say a finite collections of points $\mathcal H\subseteq\FF^n$ is an $\epsilon$-hitting set for $\mathcal F$ if for any nonzero $Q\in\mathcal F$, the evaluation $Q(\alpha)$ is nonzero for all but at most $\epsilon$-fraction of $\alpha\in\mathcal H$.
\end{result}
We need an explicit construction of $\epsilon$-hitting sets for low degree polynomials. This problem has been well studied (Bl\"aser \& Pandey 2020; Bogdanov 2005; Bshouty 2014; Cohen \& Ta-Shma 2013; Demillo \& Lipton 1978; Klivans \& Spielman 2001; Lu 2012; Schwartz 1980; Zippel 1979). For completeness, we present a construction based on sparse polynomial identity testing.
Recall that a polynomial is $s$-sparse if it has at most $s$ monomials. We need the following lemma from Agrawal et al. (2015).
\begin{result}{Lemma 2.2 (Agrawal et al. 2015, Lemma 4, restated)}
For $n,s,d\in\mathbb N^+$ and $\epsilon_0\in(0,1)$, there exist maps $w_1,w_2,\ldots,w_N:[n]\to[N\log N]$, where $N=\operatorname{poly}(n,s,\log d,\epsilon_0^{-1})$, such that for any nonzero $s$-sparse polynomial $f\in\FF[X_1,\ldots,X_n]$ of individual degree at most $d$, all but at most $\epsilon_0$-fraction of $w_i$ among $w_1,\ldots,w_N$ satisfies $f(Y^{w_i(1)},\ldots,Y^{w_i(n)})\ne0$. Moreover, the bit complexity of computing $w_1,w_2,\ldots,w_N$ is polynomial in $N$.
\end{result}
Given $n,d\in\mathbb N^+$ and $\epsilon\in(0,1)$, we construct an $\epsilon$-hitting set for $n$-variate polynomials of degree at most $d$ as follows:
\begin{enumerate}
\item Let $s=\binom{n+d}{d}$, $\epsilon_0=\epsilon/2$, and $M=\lceil\epsilon_0^{-1}dN\log N\rceil$, where $N$ is as in Lemma 2.2.
\item Let $w_1,\ldots,w_N$ be as in Lemma 2.2, which can be computed in time $\operatorname{poly}(N)$.
\item If $\operatorname{char}(\FF)=0$, let $S=[M]\subseteq\mathbb Z\subseteq\FF$. If $\operatorname{char}(\FF)=p>0$, choose the smallest $p$-power $q$ such that $q\ge M$, and choose $S$ to be a subset of $\FF_q\subseteq\FF$ of cardinality $M$. We remark that $\FF_q$ can be constructed deterministically in time $\operatorname{poly}(M,\log p)$. To see this, note that $q\le Mp$ by the minimality of $q$. If $M\le p$, then $\FF_q$ is just $\FF_p$. On the other hand, if $p<M\le q$, then $\FF_q$ can be constructed in time $\operatorname{poly}(p,[\FF_q:\FF_p])$ (see, e.g., Lenstra (1990)), which is polynomial in $M$ since $p<M$ and $q\le Mp$.
\item Finally, construct the following collection of points in $\FF^n$ of size $MN$
\[T=\{(\alpha^{w_i(1)},\ldots,\alpha^{w_i(n)}):\alpha\in S,i\in[N]\}\subseteq\FF^n.\]
\end{enumerate}
\begin{result}{Lemma 2.3}
For any nonzero polynomial $f\in\FF[X_1,\ldots,X_n]$ of degree at most $d$, we have $f(u)\ne0$ for all but at most $\epsilon$-fraction of $u\in T$. The collection $T$ has cardinality $\operatorname{poly}(\binom{n+d}{d},1/\epsilon)$ and can be computed in time $\operatorname{poly}(|T|)$.
\end{result}
\begin{proof}
Let $f\in\FF[X_1,\ldots,X_n]$ be a nonzero polynomial of degree at most $d$. Note that $f$ is trivially $s$-sparse, where $s=\binom{n+d}{d}$. So by Lemma 2.2, for all but at most $\epsilon_0$-fraction of $i\in[N]$, we have $\widetilde f_i:=f(Y^{w_i(1)},\ldots,Y^{w_i(n)})\ne0$. Consider $i\in[N]$ such that $\widetilde f_i\ne0$. Note that $\widetilde f_i$ is a univariate polynomial of degree at most $dN\log N$. So it has at most $dN\log N\le\epsilon_0M$ zeros. Therefore, by the choice of $M$, we have $f(\alpha^{w_i(1)},\ldots,\alpha^{w_i(n)})=\widetilde f_i(\alpha)\ne0$ for all but at most $\epsilon_0$-fraction of $\alpha\in S$. It follows that $f(u)\ne0$ holds for all but at most $\epsilon$-fraction of $u\in T$, as claimed. The rest of the lemma follows easily from the construction.
\end{proof}
Note that the seed length required to choose a random element in $T$ is $\log|T|=O(\log\binom{n+d}{d}+\log(1/\epsilon))$, which is optimal up to a constant factor. We have made no effort to optimize the constant hidden in $O(\cdot)$. Interested readers may find the state-of-the-art result in Bl\"aser \& Pandey (2020), which achieves the optimal constant, at least when $d=n^{o(1)}$.
\section*{Original Section 2.2: Explicit Lossless Rank Condensers (journal pp.20--21)}
We need the following lemma in the context of lossless rank condensers. The construction in the lemma was given by Forbes \& Shpilka (2012) and the lemma itself follows implicitly from the analysis of Forbes et al. (2014). It was also stated explicitly as Forbes (2014, Theorem 5.4.3).
\begin{result}{Lemma 2.4 (Forbes 2014; Forbes et al. 2014)}
Let $n\in\mathbb N^+$ and $r\in[n]$. Let $\omega\in\FF^\times$ such that the multiplicative order of $\omega$ is at least $n$. Define the $r\times n$ matrix $W=(w_{i,j})_{i\in[r],j\in[n]}$ over $\FF[X]$ by
\[w_{i,j}=(\omega^{i-1}X)^{j-1}.\]
Then for every $n\times r$ matrix $M$ over $\FF$ of rank $r$, the polynomial $\det(WM)\in\FF[X]$ is nonzero and has degree at most $r(n-r)$ after dividing out powers of $X$.
\end{result}
\begin{result}{Corollary 2.5}
Let $n,r,W$ be as in Lemma 2.4 and $\epsilon\in(0,1)$. Let $S\subseteq\FF^\times$ be a finite set of cardinality at least $r(n-r)/\epsilon$. For every $n\times r$ matrix $M$ over $\FF$ of rank $r$, we have $\operatorname{rank}(W(\alpha)M)=r$ for all but at most $\epsilon$-fraction of $\alpha\in S$, where $W(\alpha)$ denotes the matrix $(w_{i,j}(\alpha))_{i\in[r],j\in[n]}$ over $\FF$.
\end{result}
Corollary 2.5 states that the collection $\{W(\alpha):\alpha\in S\}$ of matrices is a (weak) $(r,\epsilon|S|)$-lossless rank condenser, as defined in Forbes \& Guruswami (2015). Note that for each $\alpha\in S$, we have $\operatorname{rank}(W(\alpha))=r$ and hence $W(\alpha)$ corresponds to an $(r-1)$-subspace $U_{W(\alpha)}$ of $\PP^{n-1}$. As explained in Introduction, the collection $\mathcal H=\{U_{W(\alpha)}:\alpha\in S\}$ is an $(\mathcal F,\epsilon)$-evasive $(r-1)$-subspace family on $\PP^{n-1}$, where $\mathcal F$ is the family of $(n-r-1)$-subspaces of $\PP^{n-1}$. Choosing $S$ of size $r(n-r)+1$ and $\epsilon=1-\frac1{r(n-r)+1}$ shows that the lower bound in Theorem 1.10 is achieved when $d=1$.
